% Tapis de Quillen, with Quillen's offprint and preprint (cote 162-5, typed part): archive cote 162-5, site page 10 (PDF page 11).
% typescript. Draft compared: none.
% AI-transcribed by Claude Opus 5.5 (claude-opus-5-5, adaptive thinking, effort high), Message Batches API, batch t162_full (msgbatch_01UAh3zyWPd454AeDkhnmQ33), 2026-11-13; not yet checked by eye.
\scanpage{162-5/10}
\mspage{9}

\unsure{Ces|Cas}\xout{x} bicomplexes s'identifient aux complexes de chhines  de \unsure{longueur} 2.
Donc la catégorie de (2-Cat)-groupes abéliens est équivelente à celle de
complexes de chaines de longueur 2. L'explicitation de la 2-catégorie
associée à un tel complexe de chaines est mon interprétation chérie. Il
faudrait encore expliciter \xout{lesxmorphismes}  2-Grab-catégorie  des
homomorphismes (en un sens à préciser) entre deux\xout{x} tel\xout{l}es  2-Grab-catég
% NOTE: the right edge of the page is cut; the line ends are not visible
ries strictes, en termes d'un complexe RHom.

L'idée qui se dégage de cette approche, et que Quillen on espère v
expl\unsure{i}citer, c'est que la \xout{notionxdexcatégoriexde} notion\xout{s} de groupe \uline{abélie}
dans (n-Cat) correspond à celle de complexe de \inserted{chaines de} longueur n, en raisonnant
soit mod. isomorphisme de part et d'autre, soit mod. une notion naturell
de n-équivalence d'un côté, et dans la catégorie dérivée de l'autre. Cec
doit pouvoir se préciser en utilisant des Hom$^{\bullet}$ resp. des RHom$^{\bullet}$.

Si maintenant C est une catégorie avec un ``produit tensoriel'' assoc
tif (et éventuellement unitaire), i.e. un bifoncteur avec des isomorphis
mes d'associativité (mais sans inverse ...), on peut lui associer \unsure{sans}
doute une  Gr-catégorie universelle (tout comme à un monoïde on associe
groupe symmétrisé), d'où des invariants, qu'on pourra désigner par K$_{o}$(C)
et K$_{1}$(C). On suppose qu'on pourra de même trouver, pour tout n, une
Gr-n-catégorie universelle où envoyer C de façon compatib\xout{b}le avec $\otimes$ (\unsure{tout}
cela dans un sens à préciser), dont les invariants \hand{$\pi_n$}  seraient par déf
tion les K$_{n}$(C). \unsure{Un} cas intéressant est celui où C est une catégorie ave
produits finis (par exemple une catégorie additive) et $\otimes$ est le produit\unsure{;}
\xout{mais}\over{(}{[}Mais on notera que dans le cas d'une catégorie non seulement additi
mais abélienne, il y a une notion plus adéquate du K$_{o}$ , qui devrait corr
pondre à une notion plus adéquate de la suite des K$_{i}$ , inspirée du tapis
de Roos; sans doute c'est en termes de (``vraie'') catégories triangulées

\margin{10}
